% Intisari (abstract) -- transcribed verbatim from Isi/Intisari.doc
\chapter*{INTISARI}
\addcontentsline{toc}{chapter}{INTISARI}

\noindent\textbf{OPTIMASI BETON BERTULANG PADA STRUKTUR PORTAL RUANG},
Yohan Naftali, tahun 1999, Peminatan Program Studi Struktur, Program Studi
Teknik Sipil, Fakultas Teknik, Universitas Atma Jaya Yogyakarta.

\bigskip
\noindent
Desain struktur merupakan salah satu bagian dari keseluruhan proses
perencanaan bangunan. Semakin mahalnya harga material menyebabkan harga
struktur bangunan menjadi semakin mahal, oleh karena itu untuk menekan biaya
bangunan, perlu dicari dimensi batang sehingga memberikan harga struktur
yang paling murah (optimal), tanpa melanggar kendala-kendala yang ada. Salah
satu metoda untuk mendapatkan struktur yang optimal dapat dilakukan dengan
proses optimasi. Kasus optimasi yang terjadi dalam bidang teknik pada
umumnya berupa masalah yang tidak linier, untuk itu perlu digunakan metoda
optimasi yang dapat menyelesaikan masalah tidak linier, dari dua metoda
yaitu dengan bantuan kalkulus dan metoda numerik.

Beton bertulang merupakan material yang banyak digunakan untuk membuat
struktur bangunan karena material pembentuknya mudah didapat. Balok utama
pada struktur beton bertulang langsung ditumpu oleh kolom dan dianggap
menyatu secara kaku dengan kolom, sistem ini dikatakan sebagai sistem
portal. Portal ruang merupakan sistim portal yang dianalisa secara 3
dimensi. Portal ruang memiliki 6 perpindahan yang mungkin terjadi pada
setiap titik kumpulnya.

Harga struktur yang murah tetapi tidak melanggar kendala yang ada merupakan
fungsi sasaran yang dicari dalam tugas akhir ini. Variabel desainnya berupa
dimensi penampang beton dan tulangan baja yang digunakan. Balok dan kolom
berpenampang empat persegi panjang dan uniform sepanjang bentangnya. Hasil
tugas akhir ini berupa sebuah perangkat lunak untuk optimasi beton bertulang
pada struktur portal ruang. Kode program ditulis dengan bahasa C++ dan
dikompilasi menggunakan Borland C++.

Penyelesaian masalah optimasi beton bertulang pada struktur portal ruang
memakai metoda kalkulus terlalu rumit, oleh karena itu dalam tugas akhir ini
digunakan metoda optimasi polihedron fleksibel. Dalam menganalisa struktur
portal ruang digunakan metoda kekakuan. Balok dianalisa sebagai balok
berpenampang empat persegi panjang dan bertulangan rangkap, kolom dianalisa
sebagai kolom biaksial bertulangan simetri.

Hasil yang diperoleh dengan menggunakan program optimasi beton bertulang
pada struktur portal ruang yang dikembangkan dalam tugas akhir ini dapat
mengurangi harga struktur sekitar 9,9~\% dibandingkan hitungan struktur yang
dikerjakan tanpa dioptimasi. Dalam metoda polihedron fleksibel, semakin
banyak titik coba yang digunakan, harga struktur yang didapat menjadi
semakin murah, akan tetapi membutuhkan waktu proses optimasi yang semakin
lama. Pengembangan optimasi beton bertulang pada struktur portal ruang
merupakan hal yang masih luas dan menarik untuk dipelajari, terutama masalah
pendetailan antara sambungan balok kolom dan analisis struktur dinamis yang
belum dibahas dalam tugas akhir ini.

\clearpage
